\section{Evaluation}\label{evaluation}

\subsection{Experimentelle Konfiguration}

Alle Experimente wurden mit Python~3.12 und scikit-learn auf dem CICIoT2023-Datensatz
durchgefuehrt~\cite{neto2023ciciot2023}.
Der bereinigte Datensatz umfasst 838.602 Instanzen nach Entfernung von 3.078
BACKDOOR\_MALWARE-Eintraegen (670.881 Training, 167.721 Test, Random-Seed~42).
Die Ergebnisse werden in zwei Schritten berichtet: zunaechst mit Standardhyperparametern,
anschliessend mit den durch Grid Search optimierten Parametern.

\textbf{Grid-Search-Konfiguration.}
Die Hyperparameteroptimierung wurde mit \texttt{GridSearchCV} (cv\,=\,3,
Zielfunktion: \ba{}) auf einem stratifizierten Teilmengedatensatz von 8.045 Instanzen
(5 Dateien) durchgefuehrt, um die Rechenzeit zu begrenzen.
Die durchsuchte Grille umfasst: n\_estimators $\in \{100, 200\}$,
max\_depth $\in \{3, 5\}$, learning\_rate $\in \{0{,}05, 0{,}1\}$,
subsample $\in \{0{,}8, 1{,}0\}$; insgesamt 16 Kombinationen je Stufe und Pipeline.

\subsection{Ergebnisse Pipeline A: mit PCA (39 auf 16 Merkmale)}

Die binaere Klassifikation in Stufe~1 erzielt eine \ba{} von 0{,}9971 und einen
makro-gemittelten F1-Score von 0{,}9969.
Die Mehrklassenklassifikation in Stufe~2 erzielt eine \ba{} von 0{,}5480 und einen
makro-gemittelten F1-Score von 0{,}5784.
Die Trainingsdauer betraegt 2.240 Sekunden.

\subsection{Ergebnisse Pipeline B: ohne PCA (39 Originalmerkmale)}

Die binaere Klassifikation in Stufe~1 erzielt eine \ba{} von 0{,}9993 und einen
makro-gemittelten F1-Score von 0{,}9993.
Die Mehrklassenklassifikation in Stufe~2 erzielt eine \ba{} von 0{,}6301 und einen
makro-gemittelten F1-Score von 0{,}6670.
Die Trainingsdauer betraegt 1.657 Sekunden.

\subsection{Vergleich Pipeline A und Pipeline B}

Tabelle~\ref{tab:ergebnisse} fasst die Ergebnisse beider Pipelines zusammen.

\begin{table}[h]
\centering
\caption{Balanced Accuracy und F1-Score (makro): drei Konfigurationen}
\label{tab:ergebnisse}
\begin{tabular}{llcccc}
\toprule
Pipeline & Konfig. & BA$_1$ & F1$_1$ & BA$_2$ & F1$_2$ \\
\midrule
A (PCA) & Standard & 0{,}9971 & 0{,}9969 & 0{,}5480 & 0{,}5784 \\
A (PCA) & GS-Subset & 0{,}9969 & n.\,v. & 0{,}5384 & n.\,v. \\
A (PCA) & Best HP & 0{,}9985 & 0{,}9985 & 0{,}5364 & 0{,}5644 \\
\midrule
B (39f) & Standard & 0{,}9993 & 0{,}9993 & 0{,}6301 & 0{,}6670 \\
B (39f) & GS-Subset & 0{,}9994 & n.\,v. & 0{,}6479 & n.\,v. \\
B (39f) & Best HP & 0{,}9992 & 0{,}9993 & 0{,}6285 & 0{,}6652 \\
B (39f) & Balanced & 0{,}9992 & 0{,}9993 & \textbf{0{,}7261} & 0{,}6570 \\
\bottomrule
\end{tabular}
\smallskip\newline
{\footnotesize BA$_1$, BA$_2$: Balanced Accuracy Stufe~1/2. F1: makro-gemittelt.
Standard: Standardhyperparameter auf 838.602~Instanzen. GS-Subset:
Grid-Search-Optimum auf 8.045~Instanzen. Best HP: optimierte Hyperparameter auf
838.602~Instanzen. Balanced: Best HP plus klassengewichtete Stichproben
(\texttt{sample\_weight} balanced) in Stufe~2.}
\end{table}

Tabelle~\ref{tab:ergebnisse} zeigt die Ergebnisse in vier Konfigurationen.
Pipeline~B uebertrifft Pipeline~A in allen Konfigurationen und mit allen Hyperparametern.
Das Reentrainement auf dem vollstaendigen Datensatz liefert fuer beide Pipelines eine
geringfuegig niedrigere Stufe-2-\ba{} als die Standardparameter:
Pipeline~A sinkt von 0{,}5480 auf 0{,}5364~($-$1{,}16~Punkte),
Pipeline~B sinkt von 0{,}6301 auf 0{,}6285~($-$0{,}16~Punkte).
Dieses Ergebnis belegt, dass die auf dem Teilmengedatensatz (8.045~Instanzen) gefundenen
Hyperparameter auf dem vollstaendigen Datensatz nicht besser generalisieren; dies ist
ein wichtiger methodischer Befund fuer zukunftige Studien~\cite{hosseini2025imbalance}.

\textbf{Klassengewichtung.}
Da \gb{} keinen \texttt{class\_weight}-Parameter besitzt, werden die Trainingsinstanzen
der Stufe~2 ueber \texttt{sample\_weight} invers zur Klassenhaeufigkeit gewichtet
(\glqq{}balanced\grqq{}).
Diese Massnahme hebt die \ba{} von Stufe~2 in Pipeline~B von 0{,}6285 auf
\textbf{0{,}7261}~($+$9{,}76~Punkte) und reduziert die Diskrepanz zwischen Accuracy
(0{,}7534) und \ba{} (0{,}7261) auf 0{,}0273.
Die Klassengewichtung ist damit deutlich wirksamer als die Hyperparameteroptimierung.
Sie adressiert das Klassenungleichgewicht direkt auf der Ebene der Verlustfunktion und
verbessert die Erkennung der seltenen Angriffsklassen erheblich~\cite{hosseini2025imbalance}
(siehe Tabelle~\ref{tab:perclass_b}).

Die optimierten Hyperparameter der Grid Search sind in Tabelle~\ref{tab:best_hp}
zusammengefasst.

\begin{table}[h]
\centering
\caption{Optimierte Hyperparameter aus Grid Search}
\label{tab:best_hp}
\begin{tabular}{llllll}
\toprule
Pipeline & Stufe & lr & depth & n\_est & sub \\
\midrule
A & Stufe~1 & 0{,}10 & 5 & 200 & 0{,}8 \\
A & Stufe~2 & 0{,}05 & 3 & 100 & 1{,}0 \\
B & Stufe~1 & 0{,}05 & 5 & 100 & 1{,}0 \\
B & Stufe~2 & 0{,}05 & 3 & 200 & 0{,}8 \\
\bottomrule
\end{tabular}
\smallskip\newline
{\footnotesize lr: learning\_rate; depth: max\_depth; n\_est: n\_estimators; sub: subsample.}
\end{table}

\subsection{Interpretation: Ablation Study}

Die starke Leistung von Stufe~1 in beiden Pipelines (BA jeweils ueber 0{,}997) bestaetigt,
dass der zweistufige Ansatz das Klassenungleichgewicht strukturell effektiv adressiert.
Das Ergebnis stimmt mit den Befunden von Raturi et al.\ \cite{raturi2026exploratory} ueberein,
wonach baumbasierte Modelle auf dem CICIoT2023-Datensatz besonders gut fuer die binaere
Trennung von DoS/DDoS-Verkehr geeignet sind.

Die \ba{} von Stufe~2 liegt in allen Konfigurationen unterhalb des in Raturi et al.\
\cite{raturi2026exploratory} berichteten Wertes (0{,}952).
Das Reentrainement von Pipeline~A mit optimierten Hyperparametern auf dem vollstaendigen
Datensatz liefert BA$_2$\,=\,0{,}5364, was geringfuegig unter dem Standard-Wert von 0{,}5480
liegt.
Dies legt nahe, dass die auf dem Teilmengedatensatz gefundenen Hyperparameter den
Klassenungleichgewichten des vollstaendigen Datensatzes weniger gut begegnen.
Die Diskussion in Kapitel~\ref{diskussion} geht auf diese Befunde ein.

Tabelle~\ref{tab:perclass_a} zeigt die Klassenmetriken fuer Pipeline~A (Reentrainement
mit optimierten Hyperparametern).
Der starke Recall-Kontrast zwischen den Klassen bestaetigt das strukturelle
Klassenungleichgewicht und erklaert die niedrige \ba{}:
Mirai erzielt $R$\,=\,0{,}993, waehrend Web-based auf $R$\,=\,0{,}065 und Brute-Force auf
$R$\,=\,0{,}263 abstuert.
Zugleich zeigt die einfache Accuracy von 0{,}753 (gegenueber BA\,=\,0{,}536) eine Diskrepanz
von 0{,}217, was die Einschraenkung der Accuracy als alleinige Metrik empirisch
bestaetigt~\cite{hosseini2025imbalance}.

\begin{table}[h]
\centering
\caption{Precision, Recall, F1 pro Klasse: Stufe~2 Pipeline~A (Best HP)}
\label{tab:perclass_a}
\begin{tabular}{lccc}
\toprule
Klasse & P & R & F1 \\
\midrule
Benign       & 0{,}557 & 0{,}370 & 0{,}445 \\
Brute-Force  & 0{,}850 & 0{,}263 & 0{,}402 \\
Mirai        & 1{,}000 & 0{,}993 & \textbf{0{,}997} \\
Reconnaissance & 0{,}658 & 0{,}885 & 0{,}755 \\
Spoofing     & 0{,}697 & 0{,}642 & 0{,}668 \\
Web-based    & 0{,}789 & 0{,}065 & 0{,}120 \\
\bottomrule
\end{tabular}
\smallskip\newline
{\footnotesize P: Precision, R: Recall, F1: makro-gemittelt.
Accuracy\,=\,0{,}753, BA\,=\,0{,}536 (Gap\,=\,0{,}217).}
\end{table}

Der Leistungsvorsprung von Pipeline~B gegenueber Pipeline~A bestaetigt die Ausgangshypothese:
Die \pca{}-Dimensionsreduktion beeintraechtigt die Mehrklassenklassifikation messbar.
Dieser Befund stimmt mit der Erwartung ueberein, dass die Projektion auf 16 Hauptkomponenten
klassentrennende Information verliert, die bei seltenen Angriffsklassen besonders relevant
ist.

Tabelle~\ref{tab:perclass_b} stellt den klassenweisen Recall von Pipeline~B mit Best HP
und mit Klassengewichtung gegenueber.
Die Klassengewichtung verbessert den Recall der seltenen Klassen drastisch:
Brute-Force steigt von 0{,}299 auf 0{,}646, Web-based von 0{,}192 auf 0{,}578 und
Benign von 0{,}517 auf 0{,}709.
Diese Gewinne werden durch einen Rueckgang bei Reconnaissance (0{,}907 auf 0{,}571)
erkauft, waehrend Mirai (perfekt) und Spoofing (hoch) stabil bleiben.
Der Kompromiss faellt insgesamt positiv aus, da die \ba{} als Mittel der
klassenweisen Recall-Werte von 0{,}6285 auf 0{,}7261 steigt.

\begin{table}[h]
\centering
\caption{Recall pro Klasse: Stufe~2 Pipeline~B, Best HP gegen Klassengewichtung}
\label{tab:perclass_b}
\begin{tabular}{lcc}
\toprule
Klasse & R (Best HP) & R (Balanced) \\
\midrule
Benign         & 0{,}517 & \textbf{0{,}709} \\
Brute-Force    & 0{,}299 & \textbf{0{,}646} \\
Mirai          & 1{,}000 & 1{,}000 \\
Reconnaissance & \textbf{0{,}907} & 0{,}571 \\
Spoofing       & 0{,}857 & 0{,}853 \\
Web-based      & 0{,}192 & \textbf{0{,}578} \\
\midrule
\ba{}          & 0{,}6285 & \textbf{0{,}7261} \\
\bottomrule
\end{tabular}
\smallskip\newline
{\footnotesize R: Recall (Sensitivitaet). Balanced verbessert drei seltene Klassen,
mit Einbusse bei Reconnaissance. Accuracy Balanced\,=\,0{,}7534.}
\end{table}

Abbildung~\ref{fig:cm_stage2_b} zeigt die Konfusionsmatrix von Stufe~2 fuer Pipeline~B
mit Klassengewichtung (Bewertung auf 69.773 wahren Nicht-DoS-Testinstanzen).

\begin{figure}[h]
\centering
\includegraphics[width=\columnwidth]{confusion_matrix_stage2_pipelineB}
\caption{Konfusionsmatrix Stufe~2, Pipeline~B (39 Originalmerkmale, klassengewichtet).}
\label{fig:cm_stage2_b}
\end{figure}

\subsection{SHAP-Analyse}\label{shap_analyse}

Die \shap{}-Analyse wird ausschliesslich auf den finalisierten Modellen nach dem
Reentrainement mit optimierten Hyperparametern durchgefuehrt.
Die methodische Grundlage bildet die \shap{}-Bibliothek Version~0.52 in Kombination mit
scikit-learn~1.6.

\textbf{Explainer-Auswahl.}
Fuer Stufe~1 (binaere Klassifikation) wird der \texttt{TreeExplainer} eingesetzt, der exakte
Shapley-Werte fuer baumbasierte Modelle ohne Approximation berechnet~\cite{mohale2025xai}.
Stufe~2 (Mehrklassenklassifikation) erfordert einen anderen Ansatz: der
\texttt{TreeExplainer} von \shap{} unterstuetzt \texttt{GradientBoostingClassifier} mit
mehr als zwei Klassen in Version~0.52 nicht.
Daher wird der \texttt{PermutationExplainer} mit \texttt{predict\_proba} als
Modellschnittstelle eingesetzt.
Dieser Ansatz liefert korrekte Shapley-Werte in Erwartung und ist fuer alle
Klassifikatoren anwendbar~\cite{mohale2025xai}.
Fuer Stufe~2 werden 500 Testinstanzen und 100 Hintergrundinstanzen verwendet.

\textbf{Globale Erklaerungen.}
Die globale Feature-Importance wird als Betragsmittelwert der \shap{}-Werte ueber alle
Testinstanzen berechnet ($\overline{|\phi_i|}$) und gibt an, wie stark jedes der 39
(Pipeline~B) beziehungsweise 16 Hauptkomponenten (Pipeline~A) zur Klassifikation
beitraegt.
Tabelle~\ref{tab:shap_top5} zeigt die fuenf wichtigsten Merkmale pro Angriffsklasse fuer
Pipeline~B (Originalmerkmale).
Fuer Pipeline~A entfaellt die semantische Interpretierbarkeit, da Hauptkomponenten lineare
Kombinationen der Originalmerkmale sind.

\begin{table}[h]
\centering
\caption{Top-5 SHAP-Features pro Klasse: Pipeline B (Originalmerkmale)}
\label{tab:shap_top5}
\begin{tabular}{lp{0.62\columnwidth}}
\toprule
Klasse & Top-5 Features (absteigend nach $\overline{|\phi_i|}$) \\
\midrule
Benign      & \texttt{HTTPS}, \texttt{Number}, \texttt{Max}, \texttt{syn\_flag\_number}, \texttt{Header\_Length} \\
Brute-Force & \texttt{Number}, \texttt{Min}, \texttt{Max}, \texttt{rst\_flag\_number}, \texttt{Rate} \\
Mirai       & \texttt{Number}, \texttt{Protocol\_Type}, \texttt{Min}, \texttt{rst\_flag\_number}, \texttt{syn\_flag\_number} \\
Recon.      & \texttt{rst\_flag\_number}, \texttt{Number}, \texttt{syn\_flag\_number}, \texttt{Protocol\_Type}, \texttt{Min} \\
Spoofing    & \texttt{Min}, \texttt{Number}, \texttt{Max}, \texttt{Rate}, \texttt{Protocol\_Type} \\
Web-based   & \texttt{Number}, \texttt{Protocol\_Type}, \texttt{rst\_flag\_number}, \texttt{HTTPS}, \texttt{Max} \\
\bottomrule
\end{tabular}
\smallskip\newline
{\footnotesize Stufe~1 (DoS/DDoS): \texttt{Number}, \texttt{Protocol\_Type}, \texttt{UDP}, \texttt{AVG}, \texttt{Min}.
\shap{}-Werte des klassengewichteten Pipeline-B-Modells.}
\end{table}

\textbf{Semantische Validierung.}
Die \shap{}-Feature-Importance-Ranglisten werden mit bekannten Netzwerksignaturen
verglichen (vgl.\ Tabelle~\ref{tab:shap_top5}).
Fuer Stufe~1 (DoS/DDoS) dominieren \texttt{Number} und \texttt{Protocol\_Type}, was mit
Flutangriffen und Protokollmissbrauch konsistent ist~\cite{neto2023ciciot2023}.
Die Verbindungsanzahl \texttt{Number} ist in Stufe~2 das klassenuebergreifend wichtigste
Merkmal, was die Bedeutung des Verbindungsvolumens fuer die Angriffsklassifikation
unterstreicht.
Reconnaissance und Mirai werden zusaetzlich durch die TCP-Flags \texttt{rst\_flag\_number}
und \texttt{syn\_flag\_number} erklaert, was mit Scan-Aktivitaet und SYN-basierten
Flood-Mustern konsistent ist~\cite{neto2023ciciot2023}.
Spoofing wird nicht primaer durch \texttt{ARP} oder \texttt{ICMP} erklaert, sondern durch
aggregierte Paketstatistiken (\texttt{Min}, \texttt{Number}, \texttt{Max}); dies koennte
auf Korrelationsartefakte im Datensatz hindeuten, die einer vertieften Untersuchung
beduerften.

\textbf{Ablation Study SHAP.}
Der Vergleich der \shap{}-Erklaerungen zwischen Pipeline~A und Pipeline~B zeigt den
qualitativen Effekt der \pca{}-Reduktion auf die Interpretierbarkeit.
Fuer Stufe~1 (binaere Klassifikation, Pipeline~A) identifiziert der \texttt{TreeExplainer}
PC6 ($\overline{|\phi|}$\,=\,90.171), PC1 (62.205) und PC7 (47.902) als die drei wichtigsten
Komponenten.
Fuer Stufe~2 (Mehrklassenklassifikation, Pipeline~A) dominieren PC2 und PC1 uebergreifend
ueber alle Klassen mit Abstand, gefolgt von PC7 und PC6.
Da Hauptkomponenten lineare Kombinationen der 39 Originalmerkmale sind, fehlt ihnen jede
direkte semantische Bedeutung~\cite{alharby2025ciciot}.
Ein Sicherheitsanalyst kann aus \glqq{}PC6 ist das wichtigste Feature\grqq{} kein
Angriffsverhalten ableiten.
Pipeline~B-Erklaerungen benennen konkrete Netzwerkmerkmale wie \texttt{Rate},
\texttt{syn\_flag\_number} oder \texttt{Number}, die unmittelbar mit bekannten
Angriffssignaturen vergleichbar sind.
Eine direkte Jaccard-Aehnlichkeit zwischen den Top-10-Ranglisten beider Pipelines ist
methodisch nicht sinnvoll, da die Feature-Raeume (PCA-Komponenten vs.\ Originalmerkmale)
nicht vergleichbar sind.
Die Bewertungsmetriken nach Hermosilla et al.\ \cite{hermosilla2025xai_forensic}
(Fidelitaet, Konsistenz) werden auf die Pipeline~B-Erklaerungen angewandt, da nur dort
semantisch valide Erklaerungen vorliegen.

\textbf{Lokale Erklaerungen.}
Fuer die drei bis fuenf am staerksten fehlklassifizierten Instanzen pro Pipeline wird ein
Waterfall-Plot generiert, der zeigt, welche Features die Falschklassifikation verursachen.
Diese lokale Analyse ergaenzt die globale Perspektive um individuelle Fehlermuster.

Abbildung~\ref{fig:shap_global_a} zeigt die globale \shap{}-Feature-Importance fuer
Stufe~2, Pipeline~A.
Die dominanten Hauptkomponenten (PC2, PC1, PC7, PC6) illustrieren das
Interpretierbarkeits-Problem der \pca{}: kein Feature traegt einen semantisch
verstaendlichen Namen.
Abbildung~\ref{fig:shap_global_b} zeigt im Kontrast die globale
\shap{}-Feature-Importance fuer Pipeline~B (Originalmerkmale), in der \texttt{Number}
ueber alle Klassen dominiert und konkrete Netzwerkmerkmale benannt werden.

% SHAP Global Importance -- Pipeline A (verfuegbar)
\begin{figure}[h]
\centering
\includegraphics[width=\columnwidth]{shap_global_stage2_pipelineA}
\caption{Globale \shap{}-Feature-Importance: Stufe~2, Pipeline~A (PCA).
PC-Komponenten tragen keine semantische Bedeutung.}
\label{fig:shap_global_a}
\end{figure}

\begin{figure}[h]
\centering
\includegraphics[width=\columnwidth]{shap_global_stage2_pipelineB}
\caption{Globale \shap{}-Feature-Importance: Stufe~2, Pipeline~B (klassengewichtet).
\texttt{Number} dominiert klassenuebergreifend; TCP-Flags erklaeren Scan- und Flood-Muster.}
\label{fig:shap_global_b}
\end{figure}
